\chapter{绪论}\label{chap:introduction}

\section{研究背景与意义}\label{sec:background}

\subsection{胶质母细胞瘤的临床挑战与药物发现难点}\label{sec:gbm-clinical}

胶质母细胞瘤（Glioblastoma, GBM）是成人中最常见且最具侵袭性的原发性恶性脑肿瘤。根据美国中央脑肿瘤登记处（Central Brain Tumor Registry of the United States, CBTRUS）的流行病学数据，GBM约占所有原发性脑肿瘤的14.9\%、所有胶质瘤的57.3\%，年发病率约为3.2/10万人口~\cite{Stupp2005}。诊断中位年龄为65岁，且男性发病率高于女性，比例约为1.6:1。

当前新诊断GBM的标准治疗方案，即Stupp方案，由Stupp等人于2005年确立，至今仍是GBM治疗的金标准~\cite{Stupp2005}。该方案包括最大程度安全手术切除，随后进行同步放疗（60 Gy，30次分割）和每日替莫唑胺（Temozolomide, TMZ）化疗（75 mg/m²/天），随后进行6-12个周期的TMZ辅助化疗。尽管采用了这种积极的多模式治疗策略，GBM患者的预后仍然极差：中位总生存期仅为12-15个月，五年生存率不足5\%，中位无进展生存期约为6-7个月，几乎所有患者都会出现肿瘤复发。

GBM极差的预后可归因于以下几个固有的生物学特征：

\textbf{分子异质性}：GBM在分子水平上表现出显著的肿瘤间和肿瘤内异质性。基于转录组分析，GBM可分为四种分子亚型：经典型（Classical）、间充质型（Mesenchymal）、前神经元型（Proneural）和神经型（Neural）~\cite{Verhaak2010}。经典型以表皮生长因子受体（Epidermal Growth Factor Receptor, EGFR）基因扩增和7号染色体扩增/10号染色体缺失为特征，而前神经元型常伴有血小板源性生长因子受体 α（Platelet-Derived Growth Factor Receptor Alpha, PDGFRA）基因改变以及异柠檬酸脱氢酶（Isocitrate Dehydrogenase, IDH）基因突变。这种分子异质性导致对靶向治疗的不同反应，使单靶点药物往往难以奏效~\cite{Verhaak2010}。

\textbf{血脑屏障挑战}：血脑屏障（Blood-Brain Barrier, BBB）是 GBM 治疗中最主要的生理障碍，限制了绝大多数化疗药物有效进入肿瘤区域~\cite{Pardridge2012}。BBB 是高度专业化的内皮结构，将循环血液与脑细胞外液分离，保护中枢神经系统免受潜在有害物质的侵害~\cite{Ferraris2020}。相关研究指出，一般给药系统较难穿透血脑屏障到达脑部病变部位，因此脑靶向递药系统和 BBB 穿透策略是脑部疾病药物研发中的重要方向~\cite{翁丹芳2021}。在严格的理化约束下，仅约 2\% 的小分子药物能够穿透 BBB 进入脑实质~\cite{Pardridge2012}：分子量（Molecular Weight, MW）需低于 400--500 Da，拓扑极性表面积（Topological Polar Surface Area, TPSA）需低于 90 Ų~\cite{Ertl2000}，脂水分配系数（LogP）需维持在 1--4 之间~\cite{Ferraris2020}，同时对氢键供体和受体数目也有严格限制。这些理化约束与靶点结合活性、代谢稳定性之间的平衡，使 GBM 药物设计成为典型的多目标优化问题。

\textbf{治疗耐药性}：GBM表现出对多种治疗模式的固有和获得性耐药~\cite{Fu2025}。TMZ耐药发生在约90\%的患者中，主要通过O6-甲基鸟嘌呤-DNA甲基转移酶（O6-Methylguanine-DNA Methyltransferase, MGMT）修复机制介导。血脑屏障限制了药物到达肿瘤部位，而肿瘤微环境促进免疫抑制和治疗耐药。

\subsection{AI驱动药物发现的机遇}\label{sec:ai-opportunity}

传统药物发现周期长、成本高（一款新药平均研发周期10-15年、成本超20亿美元）、失败率高。人工智能与药物发现的融合为加速新型治疗药物的发现提供了新的技术思路。已有研究指出，深度学习已广泛应用于药物靶标预测、药物-靶标相互作用预测、从头药物分子设计和ADMET性质预测等环节，为AI辅助新药研发提供了重要方法基础~\cite{戴青青2022}。近年来，几项与GBM药物发现挑战特别相关的人工智能技术取得了较快发展：

\textbf{大语言模型（Large Language Models, LLMs）}：基于Transformer的大语言模型的快速发展（如GPT-4、Claude，以及开源模型如LLaMA和Qwen2）展示了在自然语言理解、推理和生成方面前所未有的能力~\cite{Qwen2Team2024}。这些模型可以通过微调和领域特定知识整合来适应药物发现等科学领域~\cite{jin2025effective}。

\textbf{多模态学习}：整合文本、图像、图和其他数据模态的多模态大语言模型为复杂生物医学数据的统一表示提供了潜力~\cite{Liu2024GITMol}。在分子设计中，多模态模型可以联合处理分子结构（SMILES/SELFIES表示）、药理学描述和理化性质数据。

\textbf{图神经网络（Graph Neural Networks, GNNs）}：GNN已成为学习分子图表示的强大工具，能够准确预测分子性质和生成新型分子结构~\cite{Xu2019,Gilmer2017}。结合GNN的扩散模型在分子生成任务中显示出特别的前景~\cite{Xu2023}。

\textbf{多目标优化}：同时平衡多个药物性质的多参数优化（Multi-Parameter Optimization, MPO）方法变得越来越复杂~\cite{Zhou2019Exp}。这些方法可以导航目标亲和力、BBB通透性、代谢稳定性和毒性之间的复杂权衡空间。

\subsection{研究意义}\label{sec:research-scope}

GBM 发生在中枢神经系统，血脑屏障会限制许多药物进入脑部组织，候选药物在设计时就需要考虑脑部递送问题，分子的基础理化性质和潜在安全风险也会影响后续筛选。因此，GBM 药物设计是一个多目标约束问题，不能只围绕单一活性分数展开。

传统虚拟筛选方法通常需要预先构建化合物库，再基于该化合物库进行评价。针对胶质母细胞瘤疾病，这种做法存在几点明显不足。第一，化合物库的规模有限。许多可能对胶质母细胞瘤有效，尤其是具备血脑屏障穿透潜力的分子结构，并未预先包含在库中，导致筛选阶段无法触及这些分子。第二，胶质母细胞瘤具有较强的疾病特异性。EGFR 改变、VEGF/VEGFR 通路异常、IDH 突变等靶点背景会影响候选分子的设计方向。如果筛选过程只依赖通用化合物库，模型难以充分利用这些疾病信息。第三，胶质母细胞瘤药物设计需要较早考虑血脑屏障限制。传统流程通常在候选分子筛选完成后再进行理化性质、BBB 穿透性和结构风险等评价。这样容易出现对接结果较好，但不适合作为脑部候选药物的分子。

基于上述分析，本文构建 GBM-PharmaMind 框架，尝试在生成阶段引入胶质母细胞瘤疾病知识和计算评价反馈，使生成结果更贴近早期虚拟筛选需求。

在方法层面，本文探索了多模态大语言模型与 GBM 领域知识的结合方式，使模型在生成候选分子时能够参考疾病背景和药物设计要求。在应用层面，本文为 GBM 候选药物的早期虚拟筛选提供了一种计算辅助方案。生成结果可作为后续实验筛选的候选集合，从而减少早期发现中的盲目搜索。


\section{国内外研究现状}\label{sec:related-work}

过去十余年间，AI驱动的分子生成领域发展迅速，多种方法范式在从头药物设计中取得了显著进展。然而，现有方法在面向胶质母细胞瘤（GBM）这类特定疾病领域时，仍存在明显的能力不足。本节系统回顾相关研究现状，并分析现有方法与本文研究目标之间的差距。

\subsection{传统深度生成模型在分子生成中的应用}\label{sec:traditional-dl}

基于深度学习的分子生成方法经历了从无监督表示学习到条件生成、从单目标到多目标优化的演进历程。本小节综述主要方法范式及其在从头药物设计中的应用与局限性。

\textbf{变分自编码器（VAE）类模型}：VAE以无监督方式学习分子结构的连续潜在表示~\cite{GomezBombarelli2018}，能够从潜在空间采样生成新型分子。VAE学习的潜在空间具有平滑的插值特性，允许在潜在空间中插值采样以生成具有期望性质的分子。然而，VAE经常产生无效的分子结构~\cite{GomezBombarelli2018}，对特定分子特征的控制有限，且难以直接融入领域知识约束。在面向GBM等特定疾病领域时，VAE缺乏对疾病背景知识的理解能力。

\textbf{GAN与RNN类模型}：基于生成对抗网络（GAN）和循环神经网络（RNN）的方法~\cite{Segler2017}将分子生成建模为序列生成任务，直接在SMILES字符串空间进行操作。这类方法的优势在于可以借助强大的序列生成模型，但其生成的SMILES字符串可能存在无效性问题，且难以保证生成分子满足多维理化约束~\cite{Segler2017}。更重要的是，这类模型通常采用"先生成、后筛选"的二阶段范式，生成过程缺乏对物理约束的动态感知。

\textbf{Transformer-SMILES模型}：MolGPT使用Transformer-解码器架构进行分子生成~\cite{Bagal2021}，在基准数据集上展示了具有竞争力的性能。这类方法可以利用大规模预训练语言模型的能力，但同样面临SMILES有效性和多目标约束控制的挑战~\cite{Bagal2021}。MolGPT更偏向于无条件分子生成，而非针对特定靶点或疾病领域的条件生成。

\textbf{强化学习优化模型}：REINVENT 2.0结合了基于RNN的生成与使用行为克隆和基于记忆的策略更新的RL优化~\cite{Blaschke2020}。RL方法的优势在于可以直接优化复合目标，但奖励函数设计困难，且探索空间巨大可能导致训练不稳定~\cite{Zhou2019Exp}。在GBM等多目标约束场景下，设计能够同时平衡BBB穿透性、靶点选择性、毒性和合成可行性的奖励函数面临挑战。

\textbf{图神经网络与扩散模型}：图神经网络和图扩散模型在分子生成任务中取得了显著进展~\cite{Liu2024GraphDiffusion,Xu2023}。几何潜在扩散模型能够生成具有真实3D构象的分子~\cite{Xu2023}，Graph Diffusion Transformers在多条件分子生成方面展示了优异的性能~\cite{Liu2024GraphDiffusion}。然而，这些方法通常缺乏对自然语言描述和疾病领域知识的理解能力，难以直接用于"根据靶点描述生成候选分子"的任务。

传统深度生成模型的主要不足集中在三个方面：分子有效性不稳定、疾病知识利用不足，以及多目标约束控制能力有限。对于GBM任务而言，后一问题尤为突出，因为候选分子不仅需要满足靶点匹配，还要同时考虑BBB穿透倾向和基础安全性。

\subsection{大语言模型、多模态模型与平台型工具在分子生成中的应用}\label{sec:mllm-drug-discovery}

随着 Transformer 架构在自然语言处理领域的成功，大语言模型（Large Language Models, LLMs）逐渐向化学、生物医学和药物发现等科学领域扩展~\cite{jin2025effective}。与传统深度生成模型相比，LLM 能够处理靶点生物学描述、药物设计约束条件、疾病临床背景和实验文献等多模态自然语言信息~\cite{jin2025effective}，为知识密集型的药物发现任务提供了新的技术路径。

早期研究将分子表示为 SMILES 字符串，将化学语言建模视为自然语言处理任务~\cite{Bagal2021}。基于预训练语言模型的分子性质预测~\cite{jin2025effective}、分子描述生成~\cite{jin2025effective}和分子编辑~\cite{Li2025DrugPilot}等任务相继涌现，验证了大规模语言模型在化学领域的适用性。然而，这类方法依赖于 SMILES 序列的有效性~\cite{Bagal2021}，且难以精确控制生成分子的理化约束~\cite{jin2025effective}，在结构多样性和物理合理性方面存在明显不足。

针对上述问题，化学指令微调（Chemical Instruction Tuning）和领域适配技术被引入~\cite{Yu2024}。例如，Yu 等人~\cite{Yu2024}基于 Mistral-7B 构建了大规模化学指令微调数据集 LlaSMol，显著提升了模型在化学问答和性质预测任务上的表现；Jin 等人~\cite{jin2025effective}将思维链（Chain-of-Thought, CoT）推理引入分子性质预测，增强了模型的可解释性。然而，这些方法在面向特定疾病靶点生成候选分子时，仍缺乏对疾病生物学背景、靶点作用机制和药物递送约束的充分建模~\cite{jin2025effective}。

多模态分子生成模型将分子图、SMILES、蛋白口袋结构和药物文本描述映射到统一的语义表示空间~\cite{Liu2024GITMol,Li2025ChemVLM}，为跨模态条件生成提供了基础~\cite{Liu2024GITMol,Li2025ChemVLM}。GIT-Mol~\cite{Liu2024GITMol}在统一框架内处理图形、图像和文本分子科学数据；ChemVLM~\cite{Li2025ChemVLM}探索了视觉-语言模型在化学领域的零样本和少样本学习能力；DrugPilot~\cite{Li2025DrugPilot}将 LLM 与多种工具数据库集成，支持复杂的多步推理链。其中，Liu 等人~\cite{liu2025multimodal}提出的 Llamole 将大语言模型的语义理解能力与图扩散变换器的分子图生成能力相结合~\cite{liu2025multimodal}，在受控分子生成和逆合成规划任务上取得了优于 14 个基线模型的表现~\cite{liu2025multimodal}。这种图—文本协同生成架构更契合"根据靶点描述生成候选分子"的药物设计任务~\cite{liu2025multimodal}。

平台型分子设计工具通常集成分子生成、性质预测和结构优化等功能~\cite{rusinko2024aiddison,torrensfontanals2024playmolecule}，具有使用门槛低、交互便捷和部署成熟等优点~\cite{rusinko2024aiddison,torrensfontanals2024playmolecule}，适合快速探索候选分子。然而，这类工具普遍存在模型细节不透明、可解释性和可控性有限等问题~\cite{wong2025explainable}，难以灵活引入特定疾病的多模态知识库~\cite{Li2025DrugPilot}；此外，平台型工具在生成分子时容易出现骨架收敛现象~\cite{Brown2019,varela2020skeledock}，即围绕少数已知母核进行局部修饰，化学空间探索范围有限~\cite{Brown2019}，更难与外部物理评估工具集成形成闭环反馈~\cite{Li2025DrugPilot}。

尽管大语言模型与多模态方法在分子生成任务中取得了显著进展~\cite{jin2025effective,Li2025DrugPilot}，但在以下方面仍存在不足：（1）事实可靠性不足，LLM 可能生成化学无效结构或错误信息；（2）化学有效性难以稳定保证~\cite{jin2025effective}；（3）物理约束难以精确控制~\cite{jin2025effective}；（4）在 GBM 等特定疾病场景下，还需同时满足靶点匹配能力、BBB 穿透倾向和基础成药性的多目标约束~\cite{Pardridge2012}。上述挑战构成了本文方法研究的核心动机。

\subsection{GBM靶向药物发现与CNS成药性挑战}\label{sec:gbm-drug-development}

GBM 药物发现面临的核心挑战，使其与其他肿瘤药物设计存在明显区别。主要体现在以下几个方面。

\textbf{EGFR 是 GBM 中研究较多的靶点。} EGFR及其突变体 EGFRvIII 在 GBM 中较为常见，已有较多 EGFR 抑制剂结构和母核信息被公开。多种 EGFR 酪氨酸激酶抑制剂也曾用于 GBM 相关临床研究。然而，由于血脑屏障限制和旁路信号激活等因素，EGFR 靶向疗法的效果仍然有限。这提示 EGFR 候选分子设计不能只关注靶点抑制作用，还需要考虑中枢神经系统成药性。

\textbf{VEGF/VEGFR 通路与 GBM 相关血管生成过程密切相关。} GBM 的生长和侵袭与血管生成密切相关，因此 VEGF/VEGFR 通路也是重要研究方向。抗 VEGF 治疗在复发性 GBM 中已有应用，但在总体生存获益方面仍存在局限。这说明，仅抑制血管生成往往不足以控制 GBM 进展。对于 VEGF/VEGFR 靶点，候选分子设计还需要面对结合机制复杂、疗效持续性不足等问题。

\textbf{IDH1/IDH2 突变是 GBM 的重要分子标志。} IDH 突变会导致异常代谢产物积累，并影响 GBM 的分子分型。相较于部分通路靶点，IDH1 突变体的结构口袋较为明确，因此更适合用于对接评价和候选分子筛选。但 IDH1/IDH2 突变并不适用于所有 GBM 患者，所以较高的计算评分不能直接说明其具有广泛临床适用性。

\textbf{血脑屏障是 GBM 药物设计中的主要限制。} GBM 位于中枢神经系统，候选药物需要进入脑部病灶区域才可能发挥作用。血脑屏障会阻止许多外源分子进入脑组织，使不少体外活性较好的分子难以到达有效部位。因此，GBM 候选分子在设计时需要考虑分子量、极性表面积、脂溶性等理化性质。实际穿透能力还会受到主动转运、外排泵和代谢稳定性等因素影响，简单描述符难以完全反映真实情况。

\textbf{GBM 的异质性会削弱单一靶点策略的效果。} GBM 在分子水平上存在明显差异，不同亚型和突变背景会影响治疗反应。单一靶点药物往往难以覆盖全部患者，也容易受到耐药机制影响。因此，GBM 候选分子设计不能只优化某一个指标，还需要综合考虑靶点作用、BBB 相关性质和基础安全性。

综上，GBM 药物发现不仅是普通肿瘤靶向药物设计的延伸。其特殊性来自疾病本身的分子背景，也来自中枢神经系统药物递送限制。如果生成模型缺少这些信息，后续筛选中容易出现大量价值有限的候选分子。这也是本文引入 GBM 领域知识、分阶段生成策略和计算评价反馈的原因。


\section{研究内容}\label{sec:research-content}

本文围绕"如何利用多模态大语言模型生成具有GBM治疗潜力的候选药物分子"展开。研究内容主要包括三个方面：

首先，构建面向GBM候选分子生成的多模态知识输入机制，将靶点信息、临床治疗经验、参考分子和结构约束转化为模型可调用的信息。这一部分需要回答的问题是：如何有效融合GBM病理背景、药理特性、分子结构与生物活性信息，建立统一的多模态表征空间。

其次，基于Qwen2-7B-Instruct与Llamole~\cite{liu2025multimodal}构建GBM-PharmaMind框架，并结合LoRA适配~\cite{hu2022lora}和CoT-ToT推理~\cite{Wei2022,Yao2023}实现候选分子生成。这一部分需要回答的问题是：如何引导多模态大语言模型根据复杂的多目标约束（BBB穿透性、靶点选择性、药代动力学性质）进行可控的分子生成。

最后，利用AutoDock Vina~\cite{Trott2010}、BBB经验预测、DILI/hERG结构警报~\cite{Aronov2008}和MPO Reward建立生成结果的计算评价与反馈机制。这一部分需要回答的问题是：如何建立生成分子的快速、可靠虚拟验证体系，实现生成与评估的闭环反馈。


\section{本文所做工作}\label{sec:research-work}

本文围绕胶质母细胞瘤候选药物分子的早期计算生成问题展开研究。GBM 药物发现与普通分子生成任务不同，模型既要理解疾病相关靶点信息，也要考虑中枢神经系统药物设计中的血脑屏障限制。传统虚拟筛选方法通常依赖已有化合物库，筛选空间容易受到限制。通用分子生成模型虽然可以产生新的候选结构，但如果缺少 GBM 疾病背景和后续评价反馈，生成结果可能不符合实际筛选要求。基于这一问题，本文构建了 GBM-PharmaMind 候选药物分子生成与优化框架。

本文所做工作和创新点主要包括以下四个方面。

（1）构建了面向 GBM 候选分子生成的领域知识输入机制。该部分的创新性在于，将GBM靶点背景、CNS药物设计规则和结构风险提示提前写入生成输入，而不是在候选分子生成后再统一筛除。本文整理了 GBM 相关靶点信息、临床治疗经验、参考分子结构和中枢神经系统药物设计规则，并将这些内容转化为模型可使用的提示信息。通过这种方式，模型在生成候选分子时可以参考 EGFR、VEGF/VEGFR 和 IDH1/IDH2 等靶点背景，也可以提前考虑血脑屏障相关的理化要求。

（2）首次提出GBM-PharmaMind候选药物分子生成与优化框架。该框架将多模态大语言模型、分子图—文本生成模型、LoRA领域适配和物理评估反馈连接到同一流程中，使GBM候选分子生成不再局限于通用SMILES生成。以多模态大语言模型和分子图—文本生成模型为基础，引入 LoRA 进行 GBM 领域适配，使模型更贴近胶质母细胞瘤药物设计任务。相比直接使用通用模型生成 SMILES，该方法能够将疾病信息、文本描述和分子结构生成过程联系起来，从而提高候选分子生成的针对性。

（3）设计了面向GBM候选分子生成的CoT-ToT三阶段推理策略。本文将候选分子生成拆分为骨架提议、片段组装和SMILES输出三个阶段，并在每一阶段保留多个候选方向，以减少单一路径生成带来的局部收敛问题。CoT（Chain-of-Thought）用于引导模型按照较清晰的设计步骤进行推理，ToT（Tree-of-Thoughts）用于保留多个候选方向，避免生成过程过早局限在单一路径中。该策略使候选分子的生成过程更容易被检查和调整。

（4）建立了面向GBM候选分子生成过程的计算评价与反馈机制。该机制将Vina对接、BBB预测、成药性评估和结构风险检查用于生成阶段的反馈与剪枝，并在最终阶段依据DrugCLIP和SwissADME相关结果完成候选分子排序。生成阶段的评价结果会反馈到后续生成过程中，用于减少明显不符合要求的候选结构。最终候选分子的排序则依据 DrugCLIP 和 SwissADME 测得的结果完成。实验结果表明，领域适配、分阶段推理和评估反馈设计对生成质量都有积极作用。


\section{论文结构}\label{sec:structure}

本文共分为五章。第1章介绍GBM候选药物发现的研究背景、现有方法不足和本文主要工作。第2章梳理大语言模型、分子表示、图生成模型、LoRA、CoT-ToT推理和物理评估等相关技术。第3章给出GBM-PharmaMind框架设计，包括知识库、Prompt生成器、LoRA适配、CoT-ToT生成算法和物理反馈机制。第4章展示实验设置、主实验结果、对比实验和消融实验。第5章总结本文工作，分析方法不足，并讨论后续改进方向。
